\section{Conditions at Entry, the Public Sector and Displacement}\label{app:entry}

This appendix collects three tests of conditions that could shape the racial gap in persistence: the labor market at
entry, public employment, and the father's displacement scored at two ages of the son. None of them creates the
racial opening of the IGE.\footnote{Code: \texttt{49.PaperBPolicyShocks.py} (\texttt{output/paperB\_policy\_shocks.md}).}

\paragraph{Labor-market conditions at entry.} Entry conditions do not move access to good starting jobs on average and
do not create the gap. A recession at labor-market entry lowers earnings for years \citep{kahn2010long,
oreopoulos2012short}, more so for nonwhite and less-educated workers \citep{schwandt2019unlucky}. Entry is dated at 18
plus years of schooling beyond 12, and the unemployment rate at entry is that of the state where the son's parents
lived when he was 17, with state and entry-year fixed effects (Table~\ref{tab:entry}). A percentage point of state
unemployment at entry changes the probability of holding a well-paid, fast-growing job at 25--34 by 0.004 (0.008), and
the mean rate over the first ten years by 0.004 (0.010). After a slack entry, White sons' access rises with parental
income by a further 0.009$^{**}$ (0.004) per point of unemployment and log point of parental income, and Black sons'
does not ($-0.001$, s.e.\ 0.005). Wages carry the scar: nationally, a point of unemployment at entry lowers sons' wages
at 25--34 by 0.013$^{**}$ (0.005); in the state design, Black sons' wages at 35--55 are lower by 0.035$^{*}$ (0.020) and
White sons' are not (0.007, s.e.\ 0.016; difference 0.042$^{*}$, $p=0.096$). Sons who entered markets slacker than
their state and year norm show a larger racial opening of the earnings IGE than sons who entered tight markets
(0.362$^{***}$ against 0.208$^{**}$; difference 0.154, s.e.\ 0.119; Table~\ref{tab:openinggroups}), which is not
significant; the gap opens by 0.208$^{**}$ among sons who entered tight markets.

\begin{table}[!htbp]
\centering
\caption{Labor-market conditions at entry: access to good jobs and wages (PSID, sons)}
\label{tab:entry}
\resizebox{\ifdim\width>\textwidth\textwidth\else\width\fi}{!}{%
\resizebox*{!}{\ifdim\totalheight>0.85\textheight 0.85\textheight\else\totalheight\fi}{%
\begin{tabular}{lccc}
\toprule
 & \begin{tabular}[b]{@{}c@{}}Good job\\at 25--34\end{tabular} & \begin{tabular}[b]{@{}c@{}}Log wage\\25--34\end{tabular} & \begin{tabular}[b]{@{}c@{}}Log wage\\35--55\end{tabular} \\
 & (1) & (2) & (3) \\
\midrule
\multicolumn{4}{l}{\textit{A. State unemployment rate at entry (state and entry-year FE)}} \\
Black: unemployment rate (pp) & $0.001$ & $-0.005$ & $-0.035^{*}$ \\
 & (0.005) & (0.012) & (0.020) \\
White: unemployment rate (pp) & $0.002$ & $0.003$ & $0.007$ \\
 & (0.008) & (0.010) & (0.016) \\
White $-$ Black: unemployment rate (pp) & $0.002$ & $0.007$ & $0.042^{*}$ \\
 & (0.010) & (0.016) & (0.025) \\
Black: rate $\times$ log parental income & $-0.001$ & $-0.001$ & $-0.002$ \\
 & (0.005) & (0.011) & (0.017) \\
White: rate $\times$ log parental income & $0.009^{**}$ & $0.005$ & $0.010$ \\
 & (0.004) & (0.013) & (0.022) \\
Sons (Black / White) & 1,372 / 2,116 & 1,434 / 2,160 & 849 / 1,433 \\
\addlinespace
\multicolumn{4}{l}{\textit{B. Mean state unemployment rate over the first ten years}} \\
Black: unemployment rate (pp) & $0.004$ & $-0.025$ & $-0.009$ \\
 & (0.010) & (0.018) & (0.029) \\
White: unemployment rate (pp) & $0.005$ & $0.002$ & $0.020$ \\
 & (0.011) & (0.017) & (0.017) \\
White $-$ Black: unemployment rate (pp) & $0.001$ & $0.027$ & $0.030$ \\
 & (0.015) & (0.025) & (0.033) \\
Black: rate $\times$ log parental income & $0.002$ & $0.011$ & $0.022$ \\
 & (0.006) & (0.015) & (0.027) \\
White: rate $\times$ log parental income & $0.013$ & $0.001$ & $-0.008$ \\
 & (0.008) & (0.014) & (0.027) \\
Sons (Black / White) & 1,425 / 2,169 & 1,491 / 2,214 & 895 / 1,480 \\
\addlinespace
\multicolumn{4}{l}{\textit{C. National unemployment rate at entry (linear cohort trend)}} \\
Black: unemployment rate (pp) & $0.001$ & $-0.001$ & $0.018$ \\
 & (0.004) & (0.013) & (0.014) \\
White: unemployment rate (pp) & $0.002$ & $-0.010^{*}$ & $-0.008$ \\
 & (0.005) & (0.006) & (0.012) \\
White $-$ Black: unemployment rate (pp) & $0.001$ & $-0.009$ & $-0.026$ \\
 & (0.007) & (0.014) & (0.018) \\
Black: rate $\times$ log parental income & $-0.004$ & $0.008$ & $0.036^{**}$ \\
 & (0.008) & (0.012) & (0.017) \\
White: rate $\times$ log parental income & $0.003$ & $-0.001$ & $0.014$ \\
 & (0.008) & (0.012) & (0.018) \\
Sons (Black / White) & 1,722 / 2,585 & 1,803 / 2,637 & 1,147 / 1,887 \\
\bottomrule
\end{tabular}}}
\par\smallskip
\begin{minipage}{\textwidth}\footnotesize \textit{Notes:} Son-level regressions; entry year = birth year + 18 + years of schooling beyond 12. The state rate is measured in the state of the son's parental family when he was 17 (BLS LAUS, entry years from 1976); the national rate is the BLS annual average. Regressors: the rate, the rate $\times$ centred log parental income, log parental income, years of schooling, quartic in the parent's age; panels A--B add state-at-17 and entry-year fixed effects (standard errors clustered by state), panel C a linear birth-cohort trend (clustered by entry year). Good job: the son's modal job at 25--34 is a well-paid cell with high later growth (Table~\ref{tab:access}). White $-$ Black: difference of the separate regressions. $^{***}$ $p<0.01$, $^{**}$ $p<0.05$, $^{*}$ $p<0.1$. Source: \texttt{IGE/49.PaperBPolicyShocks.py} (\texttt{output/paperB\_policy\_shocks.csv}).\end{minipage}
\end{table}


\paragraph{The public sector.} Public employment is where Black sons start more often, and it does not carry the
divergence. Public employment has been a main route of Black men into stable, well-paid jobs, with pay scales that
compress later growth. 19.1 percent of Black and 12.4 percent of White sons spend most of 25--34 in federal, state or
local government, and at equal parental income White sons are 5.9 points less likely to ($-0.059^{**}$, s.e.\ 0.025).
Government starters are more often in good jobs (17 percent of Black and 42 percent of White government starters,
against 3 and 13 percent of private starters). Among sons who start in the private sector, the racial opening of the
earnings IGE is 0.213$^{***}$ (0.058), against 0.232$^{***}$ for all sons (difference $-0.019$, s.e.\ 0.035;
Table~\ref{tab:openinggroups}); access to good jobs differs by race at equal parental income by 0.048$^{***}$ (0.014);
and within the private sector the return to schooling rises with age by 0.058$^{***}$ for White and 0.027$^{**}$ for
Black sons, a difference of 0.032$^{**}$ (0.013), the same as in all sectors.

\begin{table}[!htbp]
\centering
\caption{The opening of the racial gap by entry conditions and by sector (PSID, sons)}
\label{tab:openinggroups}
\resizebox{\ifdim\width>\textwidth\textwidth\else\width\fi}{!}{%
\resizebox*{!}{\ifdim\totalheight>0.85\textheight 0.85\textheight\else\totalheight\fi}{%
\begin{tabular}{lcc}
\toprule
 & \begin{tabular}[b]{@{}c@{}}Own earnings\end{tabular} & \begin{tabular}[b]{@{}c@{}}Family income\end{tabular} \\
 & (1) & (2) \\
\midrule
\multicolumn{3}{l}{\textit{A. By labor-market conditions at entry (sons entering from 1976 with a state at 17)}} \\
Entered in a slack market & $0.362^{***}$ & $0.244^{***}$ \\
 & (0.082) & (0.085) \\
Entered in a tight market & $0.208^{**}$ & $0.181^{**}$ \\
 & (0.086) & (0.086) \\
Slack minus tight & $0.154$ & $0.063$ \\
 & (0.119) & (0.121) \\
\addlinespace
\multicolumn{3}{l}{\textit{B. By the son's sector at 25--34}} \\
All sons & $0.232^{***}$ & $0.157^{***}$ \\
 & (0.052) & (0.055) \\
Private sector at 25--34 & $0.213^{***}$ & $0.154^{**}$ \\
 & (0.058) & (0.064) \\
Government at 25--34 & $0.215^{*}$ & $0.104$ \\
 & (0.130) & (0.124) \\
\bottomrule
\end{tabular}}}
\par\smallskip
\begin{minipage}{\textwidth}\footnotesize \textit{Notes:} Opening of the White $-$ Black gap in the log IGE: mean gap at 35--49 minus mean gap at 25--34, child-year design of Table~\ref{tab:gaprobust}, bands 25--49, clusters 1968 family. Slack entry: the state unemployment rate at entry is above what state and entry-year fixed effects predict. Sector: the son's modal sector at 25--34 (government: federal, state or local; private: private company; self-employed excluded). Slack minus tight: the groups are different sons (some share families), so its standard error is approximate. $^{***}$ $p<0.01$, $^{**}$ $p<0.05$, $^{*}$ $p<0.1$. Source: \texttt{IGE/49.PaperBPolicyShocks.py} (\texttt{output/paperB\_policy\_shocks.csv}).\end{minipage}
\end{table}


\paragraph{The father's displacement at two ages of the son.} A father figure's displacement of any kind at the son's
0--17 is associated with lower earnings of the son at 25--34, by 0.079$^{***}$ (0.029), against sons never exposed with
the same pre-event family income, parental schooling and birth year, and with a coefficient of $-0.053$ (0.043) at
35--44; the change between the two ages, 0.026 (0.037), does not differ by race (White 0.047, Black $-0.019$; difference
0.066, s.e.\ 0.090; Table~\ref{tab:fatherdispage}). For family income the association is $-0.065^{**}$ at 25--34 and
$-0.090^{**}$ at 35--44. The design pools plant closures with layoffs, spreads exposure over childhood and ends at 44;
Section~\ref{sec:stability} separates the layoff in the son's teens, whose cost appears after 40.

\begin{table}[!htbp]
\centering
\caption{A childhood shock re-scored by the son's age: the father's displacement (PSID, sons)}
\label{tab:fatherdispage}
\resizebox{\ifdim\width>\textwidth\textwidth\else\width\fi}{!}{%
\resizebox*{!}{\ifdim\totalheight>0.85\textheight 0.85\textheight\else\totalheight\fi}{%
\begin{tabular}{lccc}
\toprule
 & \begin{tabular}[b]{@{}c@{}}Displaced, son at 25--34\end{tabular} & \begin{tabular}[b]{@{}c@{}}Displaced, son at 35--44\end{tabular} & \begin{tabular}[b]{@{}c@{}}Change, 35--44 minus 25--34\end{tabular} \\
 & (1) & (2) & (3) \\
\midrule
\multicolumn{4}{l}{\textit{Son's log earnings}} \\
All sons & $-0.079^{***}$ & $-0.053$ & $0.026$ \\
 & (0.029) & (0.043) & (0.037) \\
Black sons & $-0.106^{*}$ & $-0.124$ & $-0.019$ \\
 & (0.061) & (0.080) & (0.080) \\
White sons & $-0.070^{**}$ & $-0.023$ & $0.047$ \\
 & (0.034) & (0.051) & (0.042) \\
White $-$ Black &  &  & $0.066$ \\
 &  &  & (0.090) \\
\addlinespace
\multicolumn{4}{l}{\textit{Son's log family income}} \\
All sons & $-0.065^{**}$ & $-0.090^{**}$ & $-0.025$ \\
 & (0.029) & (0.043) & (0.037) \\
Black sons & $-0.032$ & $-0.084$ & $-0.052$ \\
 & (0.058) & (0.083) & (0.088) \\
White sons & $-0.076^{**}$ & $-0.087^{*}$ & $-0.011$ \\
 & (0.035) & (0.051) & (0.040) \\
White $-$ Black &  &  & $0.040$ \\
 &  &  & (0.097) \\
\bottomrule
\end{tabular}}}
\par\smallskip
\begin{minipage}{\textwidth}\footnotesize \textit{Notes:} Sons whose father figure was first displaced (plant closure or move, layoff or firing) at the son's 0--17, against sons never exposed (pseudo event at the treated's median exposure age), with a father figure working before the (pseudo) event. One stacked regression per row in which every regressor (displaced, log pre-event family income, parental schooling, birth-year fixed effects) is interacted with the son's age band; outcome: log of the son's mean earnings or family income over the band ($\geq 2$ years). Clusters 1968 family. A selection-on-observables comparison. $^{***}$ $p<0.01$, $^{**}$ $p<0.05$, $^{*}$ $p<0.1$. Source: \texttt{49.PaperBPolicyShocks.py} (\texttt{output/paperB\_policy\_shocks.csv}).\end{minipage}
\end{table}

